\begin{figure}[H]
\centering
\caption{Famílias de modelos computacionais descritas nesta seção.}
\label{fig:familias-modelos}
\setlength{\unitlength}{1mm}
\resizebox{\textwidth}{!}{%
\begin{picture}(155,92)
\small
\put(54.5,82){\framebox(46,10){Modelos computacionais}}
\put(77.5,77){\line(0,1){5}}
\put(17,77){\line(1,0){120}}
\put(17,77){\vector(0,-1){7}}
\put(54.5,77){\vector(0,-1){7}}
\put(95.5,77){\vector(0,-1){7}}
\put(135,77){\vector(0,-1){7}}
\put(0,51){\framebox(34,19){Supervisionados}}
\put(38,51){\framebox(33,19){\shortstack{Não supervisionados\\k-means, PCA}}}
\put(75,51){\framebox(41,19){\shortstack{Sequência e decisão\\CRF, aprendizado por\\reforço, ACO}}}
\put(120,51){\framebox(30,19){\shortstack{Híbridos\\ANFIS}}}
\put(17,47){\line(0,1){4}}
\put(17,47){\line(1,0){118}}
\put(18,47){\vector(0,-1){3}}
\put(57,47){\vector(0,-1){3}}
\put(96,47){\vector(0,-1){3}}
\put(135,47){\vector(0,-1){3}}
\put(0,18){\framebox(36,26){\shortstack{Lineares e\\probabilísticos\\(regressão logística,\\Naive Bayes)}}}
\put(39,18){\framebox(36,26){\shortstack{Vizinhança e\\margem\\(KNN, SVM)}}}
\put(78,18){\framebox(36,26){\shortstack{Árvores e\\combinações\\(DT, JRip, RF,\\boosting)}}}
\put(117,18){\framebox(36,26){\shortstack{Redes neurais\\(MLP, CNN,\\RNN, BERT)}}}
\put(4,4){\makebox(58,8){mais simples e legíveis}}
\put(70,8){\vector(1,0){13}}
\put(87,4){\makebox(64,8){mais flexíveis e menos interpretáveis}}
\end{picture}
}%
\fonte{Elaborado pelo autor (2026).}
\end{figure}
